In this section we introduce the notion of timed scenario expressions
(expressions for short). Expressions allow us to construct models of
behaviours of complicated systems whose constituent parts are
distributed
systems. They provide a hierarchical way of specifying systems, where
distributed scenarios are used as the building blocks.
First, we introduce a few operations for manipulating scenarios.
%\vspace{-0.5em}
\begin{definition}
  \label{def-comp-par-seq}
Let $\xi=(\E_1,\C_1)$ and $\eta=(\E_2,\C_2)$ be two normalised scenarios
over $\Sigma_\xi$ and $\Sigma_\eta$, respectively.
 \setlist{nolistsep}
\begin{itemize}[noitemsep]
%\begin{itemize}
  \item
The \emph{concatenation} of
$\xi$ and $\eta$, denoted by $\xi\circ\eta$, is
defined if $\Sigma_\xi\cap\Sigma_\eta=\emptyset$.
If defined, $\xi\circ\eta$ is the scenario
$(\E,\C_1\cup\C_2)$ over $\Sigma_\xi\cup\Sigma_\eta$, where $\E$ is
the concatenation of $\E_1$ and $\E_2$.
The semantics of $\xi\circ\eta$, denoted by $\Sem{\xi\circ\eta}$, is
$\Sem{\xi}\circ\Sem{\eta}$, that is, every behaviour of $\xi \circ \eta$
is a concatenation of some behaviour of $\xi$ and some behaviour of
$\eta$.
\item
   The \emph{choice} between two alternatives, $\xi$ and $\eta$, is denoted by
   $\xi\BAR\eta$.
   Its semantics is defined by $\Sem{\xi\BAR\eta} = \Sem{\xi}\cup\Sem{\eta}$.
\end{itemize}
\end{definition}
\noindent
If $\xi$ and $\eta$ are two scenarios, then
$\xi\BAR\eta$ is consistent if $\xi$ or $\eta$ is consistent.
If $\xi\circ\eta$ is defined, then it is consistent
%, i.e., $\Sem{\xi\circ\eta}\neq \emptyset$
if $\xi$ and $\eta$ are both consistent.

%, and $\xi||\eta$ is consistent if the DTS $\DScenario=\{\xi,\eta\}$
%is consistent.

%It is worth mentioning that the semantics of the operator $||$ is
%in accordance with the rendezvous synchronisation semantics of a DTS
%(see \Sec{sec-dts}): the semantics of a DTS is obtained by the
%parallel composition of its components.
%----------------
%\begin{obs}
%   \label{obs-comp-par-seq}
%Let $\xi$ and $\eta$ be two consistent scenarios. Then,

%$\Sem{\xi\circ\eta} = \Sem{\xi}\circ\Sem{\eta}$ and
%%$\Sem{\xi||\eta} =
%%\Sem{\DScenario}$, where $\DScenario=\{\xi,\eta\}$ is a DTS, and
%$\Sem{\xi\BAR\eta} =
%\Sem{\xi}\cup\Sem{\eta}$.
%\end{obs}
%----------------
%\vspace{-0.5em}
\begin{definition}
   \label{def-comp-par-alt-dts}
Let $\DScenario_1=\D(\xi_1,\dots,\xi_n)$ and
$\DScenario_2=\D(\eta_1,\dots,\eta_m)$ be two distributed
scenarios over $\Sigma_{\DScenario_1}$ and $\Sigma_{\DScenario_2}$,
respectively.
 \setlist{nolistsep}
\begin{itemize}[noitemsep]
%\begin{itemize}
  \item
The \emph{concatenation} of
$\DScenario_1$ and $\DScenario_2$, denoted by
$\DScenario_1\circ\DScenario_2$, is defined if
$\Sigma_{\DScenario_1}\cap\Sigma_{\DScenario_2}=\emptyset$.
If defined, $\DScenario_1\circ\DScenario_2$ is
the DTS
$\DScenario=\D(\gamma_{11},\dots,\gamma_{1m},\dots,\gamma_{n1},\dots,\gamma_{nm})$
over
$\Sigma_{\DScenario_1}\cup\Sigma_{\DScenario_2}$, where
$\gamma_{ij}=\xi_i\circ\eta_j$ (for $1\leq i\leq n$
and $1\leq j\leq m$),
and
%The semantics of $\DScenario_1\circ\DScenario_2$, denoted by
$\Sem{\DScenario_1\circ\DScenario_2}=\Sem{\DScenario}$.
%, is $\Sem{\DScenario_1}\circ\Sem{\DScenario_2}$.
\item
  The \emph{parallel composition} of $\DScenario_1$, and
  $\DScenario_2$, denoted by $\DScenario_1||\DScenario_2$, is the DTS
  $\DScenario=\D(\xi_1,\dots,\xi_n,\eta_1,\dots,\eta_m)$
  over $\Sigma_{\DScenario_1}\cup\Sigma_{\DScenario_2}$ and
  $\Sem{\DScenario_1||\DScenario_2} =\Sem{\DScenario}$.
\item
   The \emph{choice} between two alternatives,
$\DScenario_1$ and $\DScenario_2$, is denoted by
$\DScenario_1\BAR\DScenario_2$.
Its semantics is defined by $\Sem{\DScenario_1\BAR\DScenario_2} =
 \Sem{\DScenario_1}\cup\Sem{\DScenario_2}$.
  \end{itemize}
\end{definition}
%---------------
\noindent
If $\DScenario_1$ and $\DScenario_2$ are two distributed scenarios,
then $\DScenario_1\BAR\DScenario_2$ is consistent if $\DScenario_1$ or
$\DScenario_2$ is consistent.
If $\DScenario_1\circ\DScenario_2$ is defined, then it is consistent
if $\DScenario_1$ and $\DScenario_2$ are both consistent.
It is worth pointing out that, except for the case of concatenation,
$\DScenario_1$ and $\DScenario_2$ might share some event names. In
that case, the event names in
$\Sigma_{\DScenario_1}\cap\Sigma_{\DScenario_2}$ will be synchronising
events in $\DScenario_1||\DScenario_2$. If $\DScenario_1$ and
$\DScenario_2$ are both consistent, $\DScenario_1||\DScenario_2$
might not be well-formed, and, if well-formed, might not be consistent.

For example, $\DScenario_1 = \D(\xi_1,\xi_2)$, where
 $\xi_1 = (abc, \{\TT{a}{b} \geq 2\})$ and $\xi_2 =
 (eb, \{\TT{e}{b} \leq 3\})$ is a consistent DTS, so is
$\DScenario_2 = \D(\xi_3)$, where
 $\xi_3 = (efb, \{\TT{e}{b} \geq 4\})$. But
$\DScenario_1||\DScenario_2$ is inconsistent. Observe that $b$ is a
synchronising event in $\DScenario_1$, while both $b$ and $e$ are the
synchronising events in $\DScenario_1||\DScenario_1$. Because the
constraints $\TT{e}{b} \leq 3$ in $\DScenario_1$ and $\TT{e}{b}
\geq 4$ in $\DScenario_2$ cannot be satisfied simultaneously,
$\DScenario_1||\DScenario_1$ is inconsistent.
%----------------
%\begin{obs}
%   \label{obs-comp-par}
%Let $\DScenario_1$ and $\DScenario_2$ be two consistent distributed
%scenarios. Then,
%$\Sem{\DScenario_1\circ\DScenario_2} =
%\Sem{\DScenario_1}\circ\Sem{\DScenario_2}$,
%$\Sem{\DScenario_1||\DScenario_2} =
%\Sem{\DScenario}$, where
%$\DScenario=\D(\xi_1,\dots,\xi_n,\eta_1,\dots,\eta_m)$ is a DTS,
%$\DScenario_1=\D(\xi_1,\dots,\xi_n)$ and
%$\DScenario_2=\D(\eta_1,\dots,\eta_m)$,
%and
%$\Sem{\DScenario_1\BAR\DScenario_2} =
%\Sem{\DScenario_1}\cup\Sem{\DScenario_2}$.
%\end{obs}
%--------------
%\vspace{-0.5em}
\begin{definition}
  \label{def-simple-exp}
 % $S$ is a \emph{timed scenario expression} (\emph{expression}
 % for short), if $S$ is
  A \emph{timed scenario expression} (\emph{expression}
  for short) is one of the following:
  \setlist{nolistsep}
\begin{itemize}[noitemsep]
  \item
    a DTS,
    %xsequential scenario,
  \item
    $(E_1\circ E_2)$, where $E_1$ and $E_2$ are
    expressions,
  \item
    $(E_1||E_2)$, where $E_1$ and $E_2$ are
    expressions, or
    \item
      $(E_1\BAR E_2)$, where $E_1$ and $E_2$ are  expressions.
 %     \item
 %   $E_1^*$, where $E_1$ is a scenario expression.
  \end{itemize}
  \end{definition}
%------------------
\noindent
Recall that a scenario $\xi$ can be seen as the DTS
$\D(\xi)$, so it is an expression.
%In our expressions we are going to treat a sequential
%scenario as a DTS with one component.

It should be obvious that the three operators, $\circ$, $\BAR$ and $||$,
are associative. The first one is not commutative, while the other two
are.

The following observation follows directly from the definitions.
%-----------------
%\vspace{-0.5em}
\begin{obs}
  \label{obs-dist-laws}
  Distributive laws for  expressions:
   \setlist{nolistsep}
\begin{itemize}[noitemsep]
%  \begin{itemize}
  \item
    $E_1\circ(E_2||E_3)=(E_1\circ E_2)||(E_1\circ E_3)$,
  %  \item
      $(E_1||E_2)\circ E_3=(E_1\circ E_3)||(E_2\circ E_3)$
       \item
         $E_1\circ(E_2\BAR E_3)=(E_1\circ E_2)\BAR(E_1\circ E_3)$,
    %     \item
           $(E_1\BAR E_2)\circ E_3=(E_1\circ E_3)\BAR(E_2\circ E_3)$
            \item
              $E_1||(E_2\BAR E_3)=(E_1||E_2)\BAR(E_1||E_3)$,
           %    \item
      $(E_1\BAR E_2)||E_3=(E_1||E_3)\BAR(E_2||E_3)$
    \end{itemize}
  \end{obs}
%-------------------------
%\vspace{-0.5em}
\begin{definition}
  \label{def-simple-normal-form}
  An expression, $E$, is in \emph{normal form} if $E$ is
  either a DTS $\DScenario$, or of the form
  $\DScenario_1\BAR\dots\BAR\DScenario_k$, where each $\DScenario_i$
  ($1\leq i\leq k, 1<k$) is a DTS.
\end{definition}
\noindent
Using the distributive laws an
expression can be converted into normal form. It should be obvious
that the normal form of an expression is unique.
%---------------
%\vspace{-0.5em}
\begin{obs}
  \label{obs-exp-sem}
Let $E$ be an expression. The semantics of $E$, denoted by
$\Sem{E}$, is
 \setlist{nolistsep}
\begin{itemize}[noitemsep]
%\begin{itemize}
  \item
  $\Sem{\DScenario}$, if $E$ is an
  expression with normal form $\DScenario$,
\item
  $\Sem{\DScenario_1}\cup\dots\cup\Sem{\DScenario_k}$, if $E$ is an
   expression with normal form $\DScenario_1\BAR\dots\BAR\DScenario_k$
  ($1<k$).
%\item
%  $\Sem{E_1}\cup\Sem{E_2}$, if $S$ is of the form $(E_1\BAR E_2)$, where
%  $E_1$ and $E_2$ are expressions.
\end{itemize}
\end{obs}
%---------------
%\noindent
%--------------
%\vspace{-0.5em}
\begin{definition}
  \label{def-valid}
An expression, $E$, is \emph{consistent}, if
$\Sem{E}\neq\emptyset$.
  \end{definition}
\noindent
If $E$ is a consistent expression,
we say $\Sem{E}$ is the set of behaviours that are \emph{described} or
\emph{expressed} by $E$.
%-------------------
%\vspace{-0.5em}
\begin{example}
\label{ex-race}
Consider a system in which two runners compete with each other
in a race managed by a controller \cite{Beek-2023}. The race begins with the
controller issuing a signal {\tt start}, which is received by both the
runners simultaneously.
Upon completing the run, runner $\mathtt{i}$ sends a $\mathtt{finish_i}$ signal
to the controller.
So there are two alternative outcomes, depending on
which runner finishes first.
This can be modeled by the expression $\mathtt{(C_1\BAR C_2)||R_1||R_2}$,
in which $\mathtt{C_1}$ and $\mathtt{C_2}$ correspond to the two
%alternative
outcomes, each specified by a DTS:
$\mathtt{C_1}=(\mathtt{start}~ \mathtt{finish_1}~ \mathtt{finish_2},\emptyset)$,
$\mathtt{C_2}=(\mathtt{start}~ \mathtt{finish_2}~ \mathtt{finish_1},\emptyset)$,
while $\mathtt{R_1}$ and $\mathtt{R_2}$ capture
the behaviours of the runners:
$\mathtt{R_1}=(\mathtt{start}~ \mathtt{finish_1},\emptyset)$,
$\mathtt{R_2}=(\mathtt{start}~ \mathtt{finish_2},\emptyset)$.
%as shown in \Fig{fig-race}.
Observe that %because events are instantaneous,
{\tt start} is
a synchronising event between all the components,
$\mathtt{finish_1}$ is a synchronising event between the controller
(specified by $\mathtt{(C_1\BAR C_2)}$) and $\mathtt{R_1}$,
while $\mathtt{finish_2}$ is a synchronising event between the controller
and $\mathtt{R_2}$.
  \end{example}
\noindent
In the rest of the paper our examples will not feature the operator $\circ$, as
it plays no role in the problem of realisation.
